%!TEX program = xelatex
\documentclass[11pt]{article}

\usepackage[margin=1in]{geometry}
\usepackage{amsmath, amssymb, amsthm}
\usepackage{graphicx}
\usepackage{hyperref}
\usepackage{xcolor}
\usepackage{booktabs}
\usepackage{multirow}
\usepackage{float}
\usepackage[section]{placeins}
\usepackage{caption}
\usepackage{adjustbox}

\newcommand{\safeincludegraphics}[2][]{%
  \IfFileExists{#2}{\includegraphics[#1]{#2}}{%
    \fbox{\parbox[c][0.22\textheight][c]{0.9\linewidth}{\centering Missing figure:\\ \texttt{\detokenize{#2}}}}%
  }%
}
\newcommand{\figw}{0.62\linewidth}
\newcommand{\tablesize}{\footnotesize}

\setlength{\textfloatsep}{10pt plus 2pt minus 2pt}
\setlength{\floatsep}{8pt plus 2pt minus 2pt}
\captionsetup{font=small,skip=4pt}

\newtheorem{theorem}{Theorem}
\newtheorem{assumption}{Assumption}

\setlength{\parindent}{0pt}
\setlength{\parskip}{1\baselineskip}

\title{Shared Spikes and the Conditional Optimality of Quadratic Routing in Anisotropic LLM Activations}
\author{Theory Group}
\date{\today}

\begin{document}
\maketitle

% {\color{red}
% \input{comments}
% }
% \newpage

\section{Latest Updates}

\begin{enumerate}
  \item Evaluated trained models and compared routers with static specialization, consistency, and balance metrics.

  \textbf{Conclusion:} the current best static router is still the linear running-center variant, while quadratic routers remain sensitive to balance control.

  \item  Proved the initialization behavior of the vanilla quadratic router under anisotropic residual covariance.

  \textbf{Conclusion:} the raw quadratic term induces an initialization bias, which can create early load imbalance. 

  \item  Replayed the main representation experiments on raw token activations from WoS.

  \textbf{Conclusion:} the spike is not a pure nuisance or a completely common term, and the residual semantics are not confined to a near-orthogonal low-rank tail.

  \item  Used the theory and the new activation evidence to revise the router design.

  \textbf{Conclusion:} the most plausible direction is a debiased second-order router with centered quadratic features and controlled spike suppression.
\end{enumerate}

\section{Background and Theoretical Critique}

Our earlier analysis justified quadratic routing under a clean static model of semantic geometry. Two issues remained open. The first is geometric: semantic information is not fully separated from the shared spike subspace, so hard spike removal may discard useful discriminative signal. The second is optimization: even if the static geometry favors a quadratic rule, anisotropic activations can already create load skew at initialization. 

\subsection{Initial Load Imbalance of the Raw Quadratic Router}

\paragraph{Summary of Main Results.}
This subsection identifies a structural flaw of the raw quadratic router at initialization. We prove that the combination of quadratic routing and anisotropic input features induces a logit mean shift and breaks the desired uniform load distribution. As a result, the router begins from a skewed state, which training can then amplify.

Let $x \in \mathbb{R}^D$ be the router input, and let
$$x_\perp := (I-P_C)x$$
be the representation after removing the estimated common spike subspace $C$.
Let
$$z := R^\top x_\perp \in \mathbb{R}^r$$
be the projected router feature. We define the raw quadratic router logit for expert $e \in \{1,\dots,E\}$ as
$$\ell_e(z) := a_e^\top z + q_e^\top (z \odot z),$$
where $a_e, q_e \in \mathbb{R}^r$ are the router parameters for expert $e$, and $\odot$ denotes the Hadamard product.

We investigate the behavior of the router at initialization under the following rigorous assumptions.

\paragraph{Assumptions 0 (Reasonable initialization Settings).}
\begin{enumerate}
    \item[(A1)] \textbf{Feature Distribution:} The spike-removed projected feature follows a centered Gaussian distribution:
    $$z \sim \mathcal{N}(0,\Lambda), \qquad \Lambda = \mathrm{diag}(\lambda_1,\dots,\lambda_r), \qquad \lambda_i>0.$$
    Importantly, no isotropy is assumed (i.e., the eigenvalues $\lambda_i$ may be strictly distinct).

    \item[(A2)] \textbf{Initialization Scheme:} The initialization of the router parameters is independent of the input $z$. Specifically, the quadratic weights are initialized as
    $$q_{e,i} \overset{\mathrm{i.i.d.}}{\sim} \mathcal{N}\!\left(0,\frac{\sigma_q^2}{r}\right), \qquad e=1,\dots,E,\ \ i=1,\dots,r.$$

    \item[(A3)] \textbf{Small-Initialization Regime:} The parameter initialization scale is sufficiently small such that $\max_{e} \|\ell_e(z)\|$ is locally bounded, ensuring that a first-order Taylor expansion of the softmax function is asymptotically valid in expectation. Specifically, we assume the remainder term in the expansion of the softmax probabilities is stochastically dominated by $\| \ell(z) \|_2^2$.
\end{enumerate}

Define the routing probability for expert $e$ as
$$p_e(z) := \frac{\exp(\ell_e(z))}{\sum_{j=1}^E \exp(\ell_j(z))},$$
and the expected load of expert $e$ over the data distribution by
$$L_e := \mathbb{E}_z[p_e(z)\mid \theta],$$
where $\theta$ denotes the fixed, initialized router parameters.

\paragraph{Theorem 1.}
Under Assumptions (A1) and (A2), the conditional expectation of the logit for expert $e$ is given by
$$m_e := \mathbb{E}_z[\ell_e(z)\mid \theta] = q_e^\top \lambda, \qquad \lambda := (\lambda_1,\dots,\lambda_r)^\top.$$

\paragraph{Proof.}
Since $z \sim \mathcal{N}(0,\Lambda)$, we have
$$\mathbb{E}[z]=0, \qquad \mathbb{E}[z_i^2]=\lambda_i.$$
Consequently, the linear term vanishes in expectation:
$$\mathbb{E}_z[a_e^\top z\mid \theta] = a_e^\top \mathbb{E}[z] = 0.$$
For the quadratic term, by linearity of expectation, we obtain
$$
\mathbb{E}_z[q_e^\top(z\odot z)\mid\theta]
= \sum_{i=1}^r q_{e,i}\mathbb{E}[z_i^2]
= \sum_{i=1}^r q_{e,i}\lambda_i
= q_e^\top \lambda.
$$
Summing the two expectations yields $\mathbb{E}_z[\ell_e(z)\mid\theta]=q_e^\top\lambda$.
\hfill $\square$

\paragraph{Corollary 1.}
Under Assumptions (A1) and (A2), the variance of the expected logit $m_e$ over the initialization distribution of $\theta$ is
$$\mathrm{Var}_\theta(m_e) = \frac{\sigma_q^2}{r}\|\lambda\|_2^2.$$

\paragraph{Proof.}
Given $m_e=\sum_{i=1}^r q_{e,i}\lambda_i$ and that the coordinates $q_{e,i}$ are mutually independent with variance $\sigma_q^2/r$, we have
$$
\mathrm{Var}_\theta(m_e)
= \sum_{i=1}^r \lambda_i^2 \mathrm{Var}(q_{e,i})
= \frac{\sigma_q^2}{r}\sum_{i=1}^r \lambda_i^2
= \frac{\sigma_q^2}{r}\|\lambda\|_2^2.
$$
\hfill $\square$

\paragraph{Proposition 1.}
For a fixed total residual energy defined as
$$\mathrm{tr}(\Lambda) = \sum_{i=1}^r \lambda_i = T,$$
the squared norm of the eigenvalue vector satisfies
$$\|\lambda\|_2^2 \ge \frac{T^2}{r},$$
with equality holding if and only if the covariance is perfectly isotropic, i.e.,
$$\lambda_1=\cdots=\lambda_r=\frac{T}{r}.$$
Hence, under a fixed total energy budget, the isotropic feature distribution uniquely minimizes $\mathrm{Var}_\theta(m_e)$.

\paragraph{Proof.}
Applying the Cauchy--Schwarz inequality to the vectors $(1,\dots,1)$ and $(\lambda_1,\dots,\lambda_r)$, we obtain
$$
\left(\sum_{i=1}^r \lambda_i\right)^2
\le \left(\sum_{i=1}^r 1^2\right)\left(\sum_{i=1}^r \lambda_i^2\right)
= r\|\lambda\|_2^2.
$$
Rearranging terms yields $\|\lambda\|_2^2 \ge \frac{T^2}{r}$. The condition for equality in Cauchy--Schwarz dictates that $(\lambda_1,\dots,\lambda_r)$ must be collinear with $(1,\dots,1)$, which proves the uniqueness of the isotropic minimum.

\hfill $\square$

\paragraph{Theorem 2.}
Under Assumptions (A1)--(A3), the expected load for expert $e$ satisfies
$$L_e = \frac{1}{E} + \frac{1}{E}\bigl(m_e-\bar m\bigr) + \mathcal{O}\big(\mathbb{E}_z[\|\ell(z)\|_2^2 \mid \theta]\big),$$
where $\bar m := \frac{1}{E}\sum_{j=1}^E m_j$ is the empirical mean of the expected logits across experts.

\paragraph{Proof.}
Let the softmax function be denoted by $s_e(\ell) := \frac{\exp(\ell_e)}{\sum_{j=1}^E \exp(\ell_j)}$.
A first-order multivariate Taylor expansion of $s_e$ around the origin $\ell=0$ gives
$$
s_e(\ell)
= \frac{1}{E}
+ \sum_{j=1}^E \frac{1}{E}\left(\delta_{ej}-\frac{1}{E}\right)\ell_j
+ R_e(\ell),
$$
where the remainder term is bounded by $|R_e(\ell)| \le C \|\ell\|_2^2$ for some constant $C>0$ in a neighborhood of the origin.
Substituting $\ell = \ell(z)$, we have
$$
p_e(z)
= \frac{1}{E}
+ \frac{1}{E}\left(\ell_e(z)-\frac{1}{E}\sum_{j=1}^E \ell_j(z)\right)
+ R_e(\ell(z)).
$$
Taking the conditional expectation with respect to $z$ given the parameters $\theta$, and noting that $\mathbb{E}_z[\ell_e(z)\mid\theta] = m_e$, we obtain
$$
L_e
= \frac{1}{E}
+ \frac{1}{E}\left(m_e - \bar{m}\right)
+ \mathbb{E}_z[R_e(\ell(z)) \mid \theta].
$$
Because $z$ is a Gaussian vector, its fourth moments are finite, ensuring that $\mathbb{E}_z[\|\ell(z)\|_2^2 \mid \theta]$ is well-defined and finite. Consequently, the expectation of the remainder term is bounded by $\mathcal{O}\big(\mathbb{E}_z[\|\ell(z)\|_2^2 \mid \theta]\big)$, completing the proof.
\hfill $\square$

\paragraph{Conclusion.}
The theoretical analysis demonstrates that the raw quadratic term $z\odot z$ induces a non-zero mean under anisotropic residual covariance. Consequently, even prior to any training updates (at random initialization), distinct experts naturally acquire different expected logits $m_e=q_e^\top\lambda$, which directly translates to skewed average softmax loads.
Furthermore, the dispersion of this initialization-induced load imbalance is strictly proportional to $\|\lambda\|_2^2$, a quantity uniquely minimized when the feature space is perfectly isotropic. Therefore, residual anisotropy is a sufficient mathematical mechanism for inducing initial load imbalance within raw quadratic routers.

\paragraph{Remark.}
We emphasize that the derived statements strictly govern the initial load skew at $\theta = \theta_0$. While this establishes a critical initial asymmetry, it does not deterministically imply terminal expert collapse, which remains subject to the full training dynamics and regularization terms.

\section{Hypotheses and What the Experiments Test}

\begin{assumption}[Original hypothesis: shared spike with orthogonal low-rank semantics]
There exists a fixed shared spike subspace $C \subset \mathbb{R}^D$ with $\dim(C)=r_c \ll D$. For each semantic type $k$, there exists a tail subspace $U_k \subset C^\perp$ with $\dim(U_k)=r \ll D$ such that
\[
x = c + u_k + \epsilon, \qquad c \in C, \quad u_k \in U_k.
\]
The distribution of $c$ does not depend on $k$. The tail component has mean zero and covariance $\mathrm{Cov}(u_k \mid k)=\sigma_u^2 P_{U_k}$. Different semantic tails are only weakly overlapping, so the pairwise overlap between $U_k$ and $U_{k'}$ is small when $k \neq k'$.
\end{assumption}

Under this hypothesis, Experiment 1 should recover a strong shared spike, Experiment 2 should reveal a low-rank residual once a small common spike is removed, and Experiment 3 should show near-orthogonal residual tails.

\begin{assumption}[Updated hypothesis: shared spike with distributed semantics]
There exists a fixed shared spike subspace $C \subset \mathbb{R}^D$ with $\dim(C)=r_c \ll D$. For each domain or pseudo-semantic type $k$, the token activation decomposes as
\[
x = c^{\mathrm{sh}} + s_k + r_k + \epsilon, \qquad c^{\mathrm{sh}}, s_k \in C.
\]
The shared part $c^{\mathrm{sh}}$ is approximately domain-invariant, while $s_k$ is domain-conditioned structure that still lives inside the spike space. The residual component $r_k$ survives spike suppression, remains compressible rather than extremely low-rank, and occupies a residual subspace whose overlap with other domains is reduced but not negligible.
\end{assumption}

Under this hypothesis, Experiment 1 should still show strong shared dominant directions, but centering should not remove all shared geometry. Experiment 2 should show a structured residual spectrum with a nontrivial effective rank. Experiment 3 should show decreasing overlap for deeper tails, without reaching near-orthogonality.

\begin{table}[H]
\centering
\tablesize
\begin{adjustbox}{max width=\linewidth}
\begin{tabular}{p{0.14\linewidth}p{0.25\linewidth}p{0.25\linewidth}p{0.24\linewidth}}
\toprule
Experiment & Original hypothesis predicts & Updated hypothesis predicts & Current result \\
\midrule
Exp1 & Strong shared spike & Strong shared spike, and substantial shared structure after centering & Shared spike is clear; centered overlap remains high \\
Exp2 & Low-rank residual after modest spike removal & Compressible but distributed residual & Effective rank stays large after spike removal \\
Exp3 & Near-orthogonal residual tails & Deeper tails are less aligned but still overlap & Overlap decreases, but remains non-negligible \\
\bottomrule
\end{tabular}
\end{adjustbox}
\caption{How the first three experiments test the old and new hypotheses.}
\label{tab:hypothesis-test-map}
\end{table}

The evidence is consistent with the updated hypothesis rather than the original one. The shared spike is real, but semantic variation is not confined to a small orthogonal tail.

\section{Experimental Setup}

All experiments in this section are based on raw token activations rather than pooled document representations. We use the Web of Science dataset with four domains, Medical, Psychology, CS, and biochemistry, and analyze layers 4, 18, and 32. For each layer and domain, we retain up to 2048 sampled token activations. The source token pools are large enough that the sampling cap is active in every case. Experiment 1 measures cross-domain overlap before and after centering. Experiment 2 removes a shared spike of rank $r_c$ and reports the effective rank needed to explain 90\% of the residual energy. Experiment 3 fixes $r_c=4$ and measures pairwise residual tail overlap as a function of the tail start and subspace width $r$. The exp4\_like router comparison is still unfinished, so no router-level accuracy or F1 numbers are reported here.

\begin{table}[H]
\centering
\tablesize
\begin{adjustbox}{max width=\linewidth}
\begin{tabular}{lrr}
\toprule
Domain & Valid tokens seen & Sampled tokens per layer \\
\midrule
Medical & 198327 & 2048 \\
Psychology & 158442 & 2048 \\
CS & 138485 & 2048 \\
biochemistry & 185957 & 2048 \\
\bottomrule
\end{tabular}
\end{adjustbox}
\caption{Token sampling summary for the raw-activation replay on WoS. The same sampling cap is used at layers 4, 18, and 32.}
\label{tab:token-summary}
\end{table}

\section{Experimental Results}

\subsection{Experiment 1: Shared dominant directions remain strong in raw activations}

\begin{figure}[H]
\centering
\safeincludegraphics[width=\figw]{figures/exp1_mean_cosine_layer18.png}
\caption{Cross-domain cosine similarity of raw activation means at layer 18. The off-diagonal entries remain close to one, indicating a strong shared direction even before any projection.}
\label{fig:exp1-mean-cos}
\end{figure}

\begin{figure}[H]
\centering
\safeincludegraphics[width=\figw]{figures/exp1_pairwise_overlap_centered_layer18.png}
\caption{Centered pairwise overlap curves at layer 18. The overlap remains high over a wide range of $s$, which shows that shared structure survives centering and is not reducible to a single uncentered mean direction.}
\label{fig:exp1-centered-overlap}
\end{figure}

\begin{table}[H]
\centering
\tablesize
\begin{adjustbox}{max width=\linewidth}
\begin{tabular}{lrrrr}
\toprule
Layer & Centered overlap at $s=3$ & Uncentered overlap at $s=3$ & Centered overlap at $s=512$ & Uncentered overlap at $s=512$ \\
\midrule
4  & 0.8680 & 0.8943 & 0.4593 & 0.4602 \\
18 & 0.7549 & 0.8149 & 0.5784 & 0.5790 \\
32 & 0.6744 & 0.7180 & 0.4217 & 0.4222 \\
\bottomrule
\end{tabular}
\end{adjustbox}
\caption{Average pairwise subspace overlap $\|U_A^\top U_B\|_F^2 / s$ across the six WoS domain pairs. Centering mainly changes the first few dominant directions; by $s=512$ the centered and uncentered curves are nearly identical.}
\label{tab:exp1-overlap}
\end{table}

Figure \ref{fig:exp1-mean-cos} shows that the raw activation means of the four domains are already highly aligned. Table \ref{tab:exp1-overlap} gives the same message from a subspace view. At $s=3$, the centered overlap remains large in every layer, from 0.6744 to 0.8680, and the uncentered overlap is consistently larger by 0.0263 to 0.0600. By $s=512$, this gap falls below 0.001.

Centering therefore affects mainly the first few dominant directions. It does not remove the shared structure itself, because the centered overlaps remain high. The spike is not a pure nuisance term; it contains a strong common mode, and shared geometry persists after centering.

\subsection{Experiment 2: Spike removal reveals compressible but not low-rank residual structure}

\begin{figure}[H]
\centering
\safeincludegraphics[width=\figw]{figures/exp2_scree_multidomain_rc4_layer18.png}
\caption{Residual spectra after removing a rank-4 shared spike at layer 18. The four domains show closely matched scree curves, which indicates a similar residual spectral shape across domains.}
\label{fig:exp2-scree}
\end{figure}

\begin{table}[H]
\centering
\tablesize
\begin{adjustbox}{max width=\linewidth}
\begin{tabular}{lrrrrr}
\toprule
Layer & $r_c=1$ & $r_c=4$ & $r_c=16$ & $r_c=64$ & $r_c=256$ \\
\midrule
4  & 600.0 & 620.5 & 646.2 & 688.8 & 728.5 \\
18 & 231.5 & 607.5 & 666.5 & 744.0 & 834.0 \\
32 & 699.8 & 736.5 & 784.2 & 838.5 & 879.5 \\
\bottomrule
\end{tabular}
\end{adjustbox}
\caption{Average effective rank across the four WoS domains, using the rank needed to explain 90\% of the residual energy after removing a shared spike of rank $r_c$.}
\label{tab:exp2-rank}
\end{table}

The residual spectra in Figure \ref{fig:exp2-scree} decay smoothly and remain closely matched across domains, so spike removal does not destroy the shared large-scale geometry. Table \ref{tab:exp2-rank} shows, however, that the residual is not extremely low-rank. Once $r_c \ge 4$, the effective rank rises monotonically in every layer. At layer 18, for example, it increases from 607.5 at $r_c=4$ to 834.0 at $r_c=256$.

The small value at layer 18 for $r_c=1$ is also useful. Removing only one direction leaves a concentrated residual, which suggests that a single shared direction is not enough to isolate the common spike. After a modest removal such as $r_c=4$, the residual remains anisotropic and compressible, but its discriminative variation is spread across hundreds of directions. The data therefore support a compressible but distributed residual, not a very low-rank semantic tail.

\subsection{Experiment 3: Deeper residual tails are less aligned, but not near-orthogonal}

\begin{figure}[H]
\centering
\safeincludegraphics[width=\figw]{figures/exp3_mean_overlap_vs_r_layer18.png}
\caption{Mean pairwise overlap of residual tail subspaces at layer 18 after removing a rank-4 shared spike. Larger tail starts reduce overlap, but enlarging $r$ brings shared structure back.}
\label{fig:exp3-overlap}
\end{figure}

\begin{table}[H]
\centering
\tablesize
\begin{adjustbox}{max width=\linewidth}
\begin{tabular}{lrrrr}
\toprule
Layer & $(\text{start}=4, r=64)$ & $(\text{start}=64, r=64)$ & $(\text{start}=192, r=64)$ & $(\text{start}=192, r=256)$ \\
\midrule
4  & 0.4720 & 0.1772 & 0.0671 & 0.2011 \\
18 & 0.5722 & 0.2769 & 0.1146 & 0.2878 \\
32 & 0.3941 & 0.1409 & 0.0596 & 0.1816 \\
\bottomrule
\end{tabular}
\end{adjustbox}
\caption{Mean residual-tail overlap $\|U_A^\top U_B\|_F^2 / r$ across the six WoS domain pairs after removing a rank-4 shared spike.}
\label{tab:exp3-overlap}
\end{table}

Experiment 3 tests the strongest part of the earlier hypothesis. If the residual tails were close to orthogonal, overlap should stay small once we move away from the shared prefix. Table \ref{tab:exp3-overlap} does not support that claim. Deeper tail starts do reduce overlap. At layer 18 and $r=64$, the mean overlap drops from 0.5722 at start $=4$ to 0.2769 at start $=64$, and then to 0.1146 at start $=192$.

This reduction is real, but it is not enough to support near-orthogonality. The overlap never approaches zero, and it rises again when the tail width becomes large. At layer 18 and start $=192$, it increases from 0.1146 at $r=64$ to 0.2878 at $r=256$. The residual tail is therefore more separable than the spike-dominated prefix, but it still contains substantial shared structure.

\subsection{Joint implication of the first three experiments}

The three experiments support a consistent update. Raw activations contain a strong shared spike, but removing that spike does not reveal a tiny, cleanly orthogonal semantic tail. Instead, the residual remains structured, moderately high-rank, and only partially separated. This is precisely why the router should suppress the shared spike without discarding all of it, and why the quadratic branch should be corrected for anisotropic second-moment bias rather than used in its raw form.

\section{Methodology}

\subsection{Improving the second-order router from the proof}

Section 1 shows that the raw quadratic feature $z \odot z$ has expectation $\lambda$, so each expert receives an initialization-dependent mean shift $q_e^\top \lambda$. The direct fix is to center the quadratic statistic itself.

We decompose the activation into spike and residual parts,
\[
x_C = P_C x, \qquad x_\perp = (I-P_C)x,
\]
and project the residual branch to a routing feature $z = W_\perp^\top x_\perp$. We then estimate the residual second moment $\hat{\lambda}_i \approx \mathbb{E}[z_i^2]$ from a calibration stream or a running average, define a variance-normalized coordinate
\[
\tilde z_i = \frac{z_i}{\sqrt{\hat{\lambda}_i + \varepsilon}},
\]
and replace the raw quadratic feature with the centered version
\[
\phi_2(\tilde z) = \tilde z \odot \tilde z - \mathbf{1}.
\]
Under the static model, this makes the quadratic term mean-zero coordinatewise and removes the bias term behind the initialization load skew.

The experiments also show that the spike branch should not be discarded completely. Since the spike space may still contain domain-conditioned information, we use a low-dimensional spike summary instead of a hard removal:
\[
\hat x_C = V^\top (x_C - \mu_C).
\]
A practical second-order router can then use
\[
\ell_e(x) = u_e^\top \hat x_C + a_e^\top \tilde z + b_e^\top \phi_2(\tilde z).
\]
The first term preserves useful information inside the spike space. The second term keeps linear separation in the residual branch. The third term adds second-order discrimination without the mean-shift defect proved in Section 1. In practice, this router should still use a small quadratic initialization scale or a load-balance warm-up stage.

\section{Model Evaluation}

Table \ref{tab:model-eval} reports the model-level evaluation results.

\begin{table}[H]
\centering
\tablesize
\begin{adjustbox}{max width=\linewidth}
\begin{tabular}{lrrr}
\toprule
Rank & Model & Avg score & $\Delta$ vs. baseline \\
\midrule
1 & Linear + LB warm-up & 0.4540 & +0.0179 \\
2 & Linear + running center & 0.4509 & +0.0149 \\
3 & Quadratic $r=256$ + LB warm-up & 0.4493 & +0.0133 \\
4 & Quadratic $r=64$ + LB warm-up & 0.4491 & +0.0131 \\
5 & Quadratic $r=256$, no LB & 0.4374 & +0.0013 \\
6 & Linear baseline & 0.4360 & +0.0000 \\
7 & Quadratic $r=64$, no LB & 0.4322 & -0.0038 \\
\bottomrule
\end{tabular}
\end{adjustbox}
\caption{Model-level evaluation on seven tasks from the detailed evaluation report.}
\label{tab:model-eval}
\end{table}

\section{Static Router Evaluation}

We evaluate routers on three static properties: specialization, consistency, and balance. Specialization is measured with two raw statistics, the effective dimension $D_{\mathrm{eff}}$ and a fixed-r principal-subspace overlap, denoted Overlap@r with $r=50$. The routing rule is top-1 only. For $D_{\mathrm{eff}}$, we keep only experts that receive at least 256 tokens and subsample 256 tokens per valid expert before estimating the covariance spectrum. This avoids the false low-rank effect caused by unequal sample counts. Consistency uses pseudo-labels from a fixed dense teacher rather than from the evaluated model itself. Balance uses the load coefficient of variation. We then map the specialization, consistency, and balance terms to $S_{\mathrm{rank}}, S_{\mathrm{cons}}, S_{\mathrm{bal}} \in [0,1]$ and combine them with a geometric mean,
\[
RQS = (S_{\mathrm{rank}} S_{\mathrm{cons}} S_{\mathrm{bal}})^{1/3},
\]
so that failure in any one dimension is strongly penalized.

\begin{table}[H]
\centering
\tablesize
\begin{adjustbox}{max width=\linewidth}
\begin{tabular}{p{0.23\linewidth}p{0.37\linewidth}p{0.28\linewidth}}
\toprule
Metric & Definition on WoS & Purpose \\
\midrule
$D_{\mathrm{eff}}(e)$, effective dimension of routed subspace & For tokens routed to expert $e$, compute the covariance spectrum and use $D_{\mathrm{eff}}(e)=\frac{(\sum_j \lambda_j)^2}{\sum_j \lambda_j^2}$ & Tests whether an expert captures a low-dimensional manifold rather than isotropic noise \\
$\mathrm{Overlap@}r$ with $r=50$ & Compute fixed-r principal-subspace overlap at rank 50 across valid experts & Tests whether expert subspaces are geometrically separated at a shared rank scale \\
$H(E \mid C)$, routing consistency entropy & Build pseudo-labels $C$ from last-layer representations and compute the conditional entropy of expert assignments given the pseudo-label & Tests whether tokens with the same downstream target are routed consistently \\
$CV_{\mathrm{load}}$, load coefficient of variation & Let $p_e$ be the empirical token fraction routed to expert $e$, and compute the coefficient of variation of $\{p_e\}_{e=1}^E$ & Tests whether specialization is achieved without static load collapse \\
\bottomrule
\end{tabular}
\end{adjustbox}
\caption{Static router metrics. Specialization is measured with both effective dimension and fixed-r subspace overlap.}
\label{tab:router-metrics}
\end{table}

These quantities target the main question of the report. $D_{\mathrm{eff}}$ and Overlap@r ask whether experts are both internally concentrated and mutually separated. $H(E \mid C)$ asks whether routing is stable with respect to the pseudo-semantic structure in the last-layer representation. $CV_{\mathrm{load}}$ asks whether this specialization is achieved without collapse.

\begin{table}[H]
\centering
\tablesize
\begin{adjustbox}{max width=\linewidth}
\begin{tabular}{lrrrrrr}
\toprule
Rank & Model & RQS & $D_{\mathrm{eff}} \downarrow$ & Overlap@50 $\downarrow$ & $H(E \mid C) \downarrow$ & $CV_{\mathrm{load}} \downarrow$ \\
\midrule
1 & Linear + running center & 0.6038 & 63.6444 & 0.1896 & 1.4810 & 0.7898 \\
2 & Linear + LB warm-up & 0.6003 & 73.5577 & 0.2078 & 1.6600 & 0.6975 \\
3 & Quadratic $r=256$ + LB warm-up & 0.4507 & 67.1330 & 0.2085 & 1.5716 & 1.1843 \\
4 & Quadratic $r=64$ + LB warm-up & 0.3554 & 63.1984 & 0.2333 & 1.3536 & 1.3665 \\
5 & Linear baseline & 0.0001 & 77.9475 & -- & 0.0004 & 2.6456 \\
6 & Quadratic $r=64$, no LB & 0.0001 & 80.7633 & 0.1712 & 0.0459 & 2.5972 \\
7 & Quadratic $r=256$, no LB & 0.0001 & 79.2971 & 0.2477 & 0.0907 & 2.5925 \\
\bottomrule
\end{tabular}
\end{adjustbox}
\caption{Static router evaluation on the low-rank specialization benchmark. Lower $D_{\mathrm{eff}}$, lower Overlap@50, lower $H(E \mid C)$, and lower $CV_{\mathrm{load}}$ are better. ``--'' marks an unavailable overlap statistic.}
\label{tab:router-rqs}
\end{table}

Adding overlap leaves the ranking essentially unchanged. The linear running-center model remains first with an RQS of 0.6038, and the linear model with load-balance warm-up remains a close second at 0.6003. Among quadratic routers, the best result is still the $r=256$ model with load-balance warm-up, but its RQS is lower at 0.4507 because its balance term is weaker.

The overlap statistic is useful, but it should not be read in isolation. The lowest Overlap@50 comes from the quadratic no-load-balance $r=64$ model, yet that same model has $CV_{\mathrm{load}}=2.5972$ and an RQS near zero. Low overlap without balance can therefore reflect collapse rather than good specialization.

These results sharpen the current picture. Quadratic routing can still produce good specialization scores, but without explicit balance control it does not yet deliver the best static router. The evidence therefore supports second-order routing as a promising direction, but not yet as the best deployed router.

\section{Next Iteration}

The next iteration should complete a regression experiment and report static WoS router results under the same pseudo-label protocol used in the specialization benchmark. That step should tell us whether the debiased quadratic router improves expert specialization on the same activation geometry studied in Exp1--Exp3.

After that, the report should add one focused ablation over spike rank, quadratic dimension, and second-moment centering. The goal is to identify whether the remaining gap comes from the quadratic feature itself or from the way balance is enforced.

If those static results are stable, the final step is to add a small training-dynamics appendix. That appendix should be narrow: early load skew, expert coverage, and whether the centered quadratic branch reduces collapse during the first stage of training.

\end{document}
